\documentclass[11pt]{article}
\usepackage{ap-psych}
\title{Unit 3: Development and Learning: Question Set}
\author{Paul Wang}
\date{November 2026}
\docpart{Unit 3 \textperiodcentered\ Question Set}

\begin{document}
\maketitle

\begin{abstract}
Exam-style multiple choice on all nine Unit 3 topics, with a shared-study set on attachment and a full Evidence-Based Question on discipline and children's behavior. Learning items use everyday scenarios where the trap is the sign (positive vs.\ negative) or the direction (reinforcement vs.\ punishment); development items describe a child and ask for the stage. All items and sources are original and fictional, written for practice. Allow about 32 minutes for the multiple choice and 45 minutes for the EBQ.
\end{abstract}

\scoringline{MCQ: \blank\ / \qtotal \qquad EBQ: \blank\ / 7 \qquad Target: 22+ MCQ and 5+ EBQ \qquad Pace: 72 s per MCQ}

\section{Development: themes, physical, gender (3.1--3.3)}

\begin{q}{2}{3.1 Methods}{Testing the same children repeatedly over years = longitudinal.}
To learn whether shy toddlers become shy adults, the best design is
\ch{a cross-sectional study of toddlers and adults}{a longitudinal study following the same children into adulthood}{a case study of one shy adult}{naturalistic observation of a preschool}
\end{q}

\begin{q}{4}{3.2 Newborn reflexes}{Turning toward a cheek touch and opening the mouth = rooting reflex.}
When her cheek is stroked, a newborn turns her head toward the touch and opens her mouth. This is the
\chh{grasping reflex}{Moro reflex}{sucking reflex}{rooting reflex}
\end{q}

\begin{q}{1}{3.2 Maturation}{Milestones that unfold on a biological timetable despite practice differences = maturation.}
Infants around the world begin walking at roughly the same age regardless of how much their parents practice with them. This best supports the role of
\chh{maturation}{teratogens}{scaffolding}{imprinting}
\end{q}

\begin{q}{3}{3.3 Gender}{Culturally shared expectations of behavior for each gender = gender roles.}
In one culture men are expected to cook for guests; in another, cooking for guests is seen as women's work. These differences reflect
\chh{primary sex characteristics}{biological sex}{gender roles}{sexual orientation}
\end{q}

\section{Cognitive development and language (3.4--3.5)}

\begin{q}{4}{3.4 Piaget}{Searching for a hidden toy = object permanence, end of sensorimotor stage.}
A 10-month-old pulls a blanket off a toy he watched being hidden. He has developed
\chh{conservation}{egocentrism}{theory of mind}{object permanence}
\end{q}

\begin{q}{2}{3.4 Egocentrism}{Assuming others see what you see = preoperational egocentrism.}
On a video call, a 3-year-old holds up a drawing to her own face and says, ``Look, Grandma!'' without turning it toward the camera. This shows
\chh{object permanence}{egocentrism}{accommodation}{formal operations}
\end{q}

\begin{q}{1}{3.4 Formal operations}{Systematic hypothetical reasoning = formal operational.}
Asked how to test which of three factors affects a pendulum's swing, a student changes one factor at a time while holding the others constant. She is showing
\ch{formal operational thinking}{concrete operational thinking}{preoperational thinking}{the imaginary audience}
\end{q}

\begin{q}{3}{3.4 Theory of mind}{Predicting where Sally will look based on her false belief = theory of mind.}
A child correctly predicts that a puppet will look for a marble where the puppet left it, not where the marble was moved while the puppet was away. The child has
\ch{conservation}{object permanence}{a theory of mind}{a critical period}
\end{q}

\begin{q}{2}{3.4 Accommodation}{Changing a schema to fit new information = accommodation.}
After learning that whales are mammals, a student revises her belief that ``everything that lives in the ocean is a fish.'' She has
\chh{assimilated}{accommodated}{overgeneralized}{conserved}
\end{q}

\begin{q}{4}{3.5 Language}{Content words only, ``want cookie'' = telegraphic (two-word) speech.}
A 2-year-old says ``Daddy go'' and ``more juice.'' This is
\chh{babbling}{overgeneralization}{the one-word stage}{telegraphic speech}
\end{q}

\begin{q}{1}{3.5 Critical period}{Late first exposure impairs grammar mastery = critical period.}
Deaf adults who were first exposed to sign language at age 14 show less fluent grammar than those exposed from birth. This supports the idea of
\chh{a critical period}{linguistic relativity}{a teratogen}{assimilation}
\end{q}

\section{Social-emotional development (3.6)}

\begin{stimulus}
\textbf{Study D.} Researchers observed 120 one-year-olds in the Strange Situation. Results:
\begin{center}
\begin{tabular}{lc}
\toprule
Classification & Percent of infants\\\midrule
Secure & 62\\
Insecure-avoidant & 18\\
Insecure-anxious (ambivalent) & 14\\
Disorganized & 6\\
\bottomrule
\end{tabular}
\end{center}
Infants rated securely attached at age 1 were rated by teachers as more socially competent at age 6 (\(r = .38\)).
\end{stimulus}

\begin{q}{3}{3.6 Attachment}{Little distress at departure and ignoring the parent = avoidant.}
In Study D, an infant who barely notices the parent leaving and turns away on the parent's return would be classified as
\chh{secure}{anxious}{avoidant}{disorganized}
\end{q}

\begin{q}{2}{3.6 Data}{About 74 of 120 infants were secure (62\%); a majority.}
Which statement about Study D is accurate?
\ch{Most infants were insecurely attached.}{A majority of infants were securely attached.}{Disorganized attachment was the second most common.}{The correlation proves that secure attachment causes social competence.}
\end{q}

\begin{q}{4}{3.6 Correlation}{r = .38: moderate positive relationship, not proof of causation.}
The correlation of \(r = .38\) in Study D indicates that
\ch{attachment at age 1 determines social skills at age 6}{there is a strong negative relationship}{teachers caused the attachment classification}{securely attached infants tended to be rated more socially competent later, a moderate relationship}
\end{q}

\begin{q}{1}{3.6 Erikson}{Adolescent search for self = identity vs. role confusion.}
A high-school junior tries out different friend groups, styles, and career goals while wondering who she really is. Erikson would say she is in the stage of
\ch{identity vs.\ role confusion}{intimacy vs.\ isolation}{competence vs.\ inferiority}{autonomy vs.\ shame and doubt}
\end{q}

\begin{q}{3}{3.6 Parenting}{High warmth, few demands = permissive.}
Parents who are loving but set almost no rules or limits are
\chh{authoritative}{authoritarian}{permissive}{neglectful}
\end{q}

\section{Learning (3.7--3.9)}

\begin{q}{2}{3.7 Classical}{The thunder (loud noise) naturally causes the startle: UCS.}
A dog that has heard thunder during many storms now trembles when the sky darkens. In this example, the thunder is the
\chh{conditioned stimulus}{unconditioned stimulus}{conditioned response}{neutral stimulus}
\end{q}

\begin{q}{4}{3.7 Generalization}{Fear spreading to similar stimuli = generalization.}
After being bitten by a German shepherd, a child becomes afraid of all large dogs. This is
\chh{discrimination}{extinction}{higher-order conditioning}{generalization}
\end{q}

\begin{q}{1}{3.7 Extinction}{Presenting the CS repeatedly without the UCS weakens the CR.}
A therapist has a client repeatedly enter elevators without anything bad happening until the client's panic fades. In classical-conditioning terms, the therapist is producing
\chh{extinction}{acquisition}{spontaneous recovery}{shaping}
\end{q}

\begin{q}{3}{3.8 Reinforcement}{Adding something pleasant to increase behavior = positive reinforcement.}
A manager gives employees a \$50 bonus for every month without a safety violation, and employees begin following safety rules more consistently. The bonus is
\ch{negative reinforcement}{positive punishment}{positive reinforcement}{a primary reinforcer}
\end{q}

\begin{q}{2}{3.8 Schedules}{Pay after a fixed number of responses = fixed ratio.}
A strawberry picker is paid \$5 for every full crate. She is on a
\chh{variable ratio schedule}{fixed ratio schedule}{fixed interval schedule}{variable interval schedule}
\end{q}

\begin{q}{4}{3.8 Schedules}{Reward for the first response after a set time: fixed interval; responding spikes near the deadline.}
Students who know a quiz is given every Friday study little early in the week and cram on Thursday. This pattern is typical of a
\chh{variable ratio schedule}{continuous schedule}{variable interval schedule}{fixed interval schedule}
\end{q}

\begin{q}{1}{3.8 Learned helplessness}{After uncontrollable failures, the organism stops trying = learned helplessness.}
After failing several tests no matter how much he studied, Eli stops studying even in a new class where effort would pay off. This is
\chh{learned helplessness}{instinctive drift}{latent learning}{insight}
\end{q}

\begin{q}{3}{3.8 Secondary reinforcers}{Tokens gain value only through what they buy: secondary reinforcers.}
Patients in a hospital ward earn tokens for completing chores and trade them for snacks and privileges. The tokens are
\chh{primary reinforcers}{punishers}{secondary reinforcers}{unconditioned stimuli}
\end{q}

\begin{q}{2}{3.9 Insight}{Sudden solution after a pause = insight learning (Köhler).}
After staring at a puzzle for several minutes, Lena suddenly sees the solution and solves it immediately. This is
\chh{latent learning}{insight learning}{shaping}{observational learning}
\end{q}

\section{Evidence-Based Question}

\begin{frq}{Question 2 \textperiodcentered\ EBQ \textperiodcentered\ 7 points \textperiodcentered\ 45 minutes}
\small
\textbf{Directions.} Read the three summarized sources (fictional, written for practice) and answer the prompt. Cite sources by author and year.

\textbf{Prompt.} Using evidence from at least two sources, develop and justify an argument about the most effective way for adults to reduce children's problem behavior.

\medskip
\textbf{Source 1.} Ferreira \& Odum (2022). \emph{Method}: Meta-analysis of 61 longitudinal studies (total \(N \approx 140{,}000\) children, 13 countries). \emph{Results}: Physical punishment (spanking) at one time point predicted \emph{more} aggression and rule-breaking two to five years later (average \(r = .21\)), even after controlling for earlier behavior. No study found that spanking improved long-term compliance.

\textbf{Source 2.} Nakamura (2023). \emph{Participants}: 24 second-grade classrooms in Japan. \emph{Method}: Experiment. Classrooms were randomly assigned to a reinforcement program (teachers gave specific praise and points exchangeable for class privileges when students followed rules) or to usual practice. \emph{Results}: After 10 weeks, observed disruptive behaviors fell 41\% in program classrooms and 4\% in control classrooms (\(p < .001\)).

\textbf{Source 3.} Abebe, Lund, \& Cho (2021). \emph{Participants}: 830 families in Ethiopia, Norway, and South Korea. \emph{Method}: Longitudinal; parenting style was rated from home observations at age 4 and self-regulation was tested at age 10. \emph{Results}: Children of authoritative parents scored highest on self-regulation in all three countries; children of authoritarian parents scored lowest.

\medskip
\begin{enumerate}[label=(\Alph*),leftmargin=2em,itemsep=3pt]
\item Propose a specific and defensible claim based in psychological science that responds to the prompt.
\item (i) Support your claim using at least one piece of specific and relevant evidence from one of the sources. (ii) Explain how the evidence from Part B(i) supports your claim using a psychological perspective, theory, concept, or research finding learned in AP Psychology.
\item (i) Support your claim using an additional piece of specific and relevant evidence from a different source. (ii) Explain how the evidence from Part C(i) supports your claim using a different psychological perspective, theory, concept, or research finding learned in AP Psychology.
\end{enumerate}
\end{frq}

\makekey

\subsection*{EBQ scoring guide (practice rubric, 7 points)}
\small
Modeled on the published EBQ structure (claim, two pieces of cited evidence, reasoning with course content).
\begin{compare}{Part}{Earns credit}{Common miss}
A. Claim (1) & ``Reinforcing desired behavior within warm, rule-based (authoritative) parenting reduces children's problem behavior more effectively than physical punishment.'' & ``Discipline matters''; a claim the sources cannot support (``spanking always causes violence''). \\
B(i). Evidence (1) & Nakamura (2023): random assignment to a praise-and-points program cut disruptive behavior 41\% vs.\ 4\% in control classrooms. & ``A study showed rewards work'' with no source or numbers. \\
B(ii). Reasoning (2) & 1 pt: links evidence to claim. 1 pt: course concept applied, e.g., \emph{positive reinforcement} and \emph{secondary reinforcers} (points) increase rule-following; the random assignment allows a causal conclusion. & Defining reinforcement in general without tying it to the classroom result. \\
C(i). Evidence (1) & Ferreira \& Odum (2022): spanking predicted more aggression two to five years later (\(r = .21\)) across 61 studies. \emph{Or} Abebe et al.\ (2021): children of authoritative parents had the highest self-regulation in three countries. & Second fact from the same source used in B. \\
C(ii). Reasoning (2) & 1 pt: links to claim. 1 pt: a different concept, e.g., \emph{observational learning} (children model the aggression they experience, as in Bandura's Bobo doll study); \emph{positive punishment} suppresses behavior short-term without teaching alternatives; or \emph{authoritative parenting}. A strong answer notes these data are correlational, so they support but cannot prove causation. & Saying the meta-analysis proves spanking \emph{causes} aggression. \\
\end{compare}

\end{document}
